\section*{NeurIPS Paper Checklist}

\begin{enumerate}

\item {\bf Claims}
    \item[] Question: Do the main claims made in the abstract and introduction accurately reflect the paper's contributions and scope?
    \item[] Answer: \answerYes{}
    \item[] Justification: The abstract and \S\ref{sec:intro} explicitly state the four contributions (SkelKD, Hierarchical Compressor, Dual-path Condition Network, OccluRadar-Human) and all claims are supported by experiments in \S\ref{sec:exp}.

\item {\bf Limitations}
    \item[] Question: Does the paper discuss the limitations of the work performed by the authors?
    \item[] Answer: \answerYes{}
    \item[] Justification: Limitations are discussed at the end of \S\ref{sec:conclusion}: dependence on depth-camera skeleton, limited OccluRadar-Human scale, and absence of downstream task integration.

\item {\bf Theory assumptions and proofs}
    \item[] Question: For each theoretical result, does the paper provide the full set of assumptions and a complete (and correct) proof?
    \item[] Answer: \answerNA{}
    \item[] Justification: This paper presents an empirical method; it contains no theorems or formal proofs.

\item {\bf Experimental result reproducibility}
    \item[] Question: Does the paper fully disclose all the information needed to reproduce the main experimental results of the paper to the extent that it affects the main claims and/or conclusions of the paper (regardless of whether the code and data are provided or not)?
    \item[] Answer: \answerYes{}
    \item[] Justification: \S\ref{sec:impl} provides all architecture dimensions, optimizer settings, learning rates, batch sizes, epochs, and data splits. Full code and checkpoints will be released (Appendix~\ref{app:repro}).

\item {\bf Open access to data and code}
    \item[] Question: Does the paper provide open access to the data and code, with sufficient instructions to faithfully reproduce the main experimental results, as described in supplemental material?
    \item[] Answer: \answerNo{}
    \item[] Justification: Code and checkpoints will be released upon paper acceptance (anonymous URL provided in Appendix~\ref{app:repro}). MM-Fi is already publicly available; OccluRadar-Human will be released simultaneously.

\item {\bf Experimental setting/details}
    \item[] Question: Does the paper specify all the training and test details (e.g., data splits, hyperparameters, how they were chosen, type of optimizer) necessary to understand the results?
    \item[] Answer: \answerYes{}
    \item[] Justification: \S\ref{sec:impl} lists all hyperparameters, optimizer type, learning rates, EMA decay, gradient clipping, warmup, data splits (train S01--S07, val S08--S10), and hardware.

\item {\bf Experiment statistical significance}
    \item[] Question: Does the paper report error bars suitably and correctly defined or other appropriate information about the statistical significance of the experiments?
    \item[] Answer: \answerNo{}
    \item[] Justification: Following standard practice in point cloud generation (LION, TIGER, LDT), we report single-run metrics; error bars are not reported due to the high computational cost of retraining diffusion models on a single RTX 3090.

\item {\bf Experiments compute resources}
    \item[] Question: For each experiment, does the paper provide sufficient information on the computer resources (type of compute workers, memory, time of execution) needed to reproduce the experiments?
    \item[] Answer: \answerYes{}
    \item[] Justification: \S\ref{sec:impl} specifies single NVIDIA RTX 3090 (24\,GB); Appendix~\ref{app:repro} provides wall-clock training times for each stage.

\item {\bf Code of ethics}
    \item[] Question: Does the research conducted in the paper conform, in every respect, with the NeurIPS Code of Ethics?
    \item[] Answer: \answerYes{}
    \item[] Justification: All human subject data was collected with informed consent; no private or sensitive information is exposed; the technology is for healthcare/safety applications with low misuse risk.

\item {\bf Broader impacts}
    \item[] Question: Does the paper discuss both potential positive societal impacts and negative societal impacts of the work performed?
    \item[] Answer: \answerYes{}
    \item[] Justification: Positive impacts include non-invasive human sensing for healthcare and smart environments. A potential negative impact is use for unauthorized surveillance; however, mmWave radar requires close proximity and specialized hardware, substantially limiting this risk.

\item {\bf Safeguards}
    \item[] Question: Does the paper describe safeguards that have been put in place for responsible release of data or models that have a high risk for misuse?
    \item[] Answer: \answerNA{}
    \item[] Justification: The model generates human body point clouds in controlled lab settings; it poses no high-risk misuse potential comparable to large language models or image generators.

\item {\bf Licenses for existing assets}
    \item[] Question: Are the creators or original owners of assets (e.g., code, data, models), used in the paper, properly credited and are the license and terms of use explicitly mentioned and properly respected?
    \item[] Answer: \answerYes{}
    \item[] Justification: MM-Fi~\cite{yang2023mmfi} is cited and used under its public research license. All baseline methods (PCN, PoinTr, SnowflakeNet, LAKe-Net, LION, TIGER) are cited. The LDT backbone is cited and used as a backbone; no code is redistributed.

\item {\bf New assets}
    \item[] Question: Are new assets introduced in the paper well documented and is the documentation provided alongside the assets?
    \item[] Answer: \answerYes{}
    \item[] Justification: OccluRadar-Human is described in \S\ref{sec:occludata} (sensor setup, collection protocol, data content, subject demographics). A data card and collection instructions will be released with the dataset.

\item {\bf Crowdsourcing and research with human subjects}
    \item[] Question: For crowdsourcing experiments and research with human subjects, does the paper include the full text of instructions given to participants and screenshots, if applicable, as well as details about compensation?
    \item[] Answer: \answerYes{}
    \item[] Justification: OccluRadar-Human was collected with consenting subjects. Collection instructions and consent form templates are included in the supplementary material released with the dataset.

\item {\bf Institutional review board (IRB) approvals or equivalent for research with human subjects}
    \item[] Question: Does the paper describe potential risks incurred by study participants, whether such risks were disclosed to the subjects, and whether Institutional Review Board (IRB) approvals (or an equivalent approval/review based on the requirements of your country or institution) were obtained?
    \item[] Answer: \answerYes{}
    \item[] Justification: The OccluRadar-Human data collection protocol was reviewed and approved by the authors' institutional ethics committee (details omitted for anonymity; will be disclosed in the camera-ready version). Participants were informed of all procedures; no physical risks are involved (passive mmWave sensing at $<$1\,mW transmit power).

\item {\bf Declaration of LLM usage}
    \item[] Question: Does the paper describe the usage of LLMs if it is an important, original, or non-standard component of the core methods in this research?
    \item[] Answer: \answerNA{}
    \item[] Justification: LLMs are not used as a core component of SKD-Net's methodology, training, or evaluation. LLM tools were used only for writing assistance and do not affect scientific content.

\end{enumerate}
